% TikZ code for Figure 7: Overall Fairness-Aware Histopathology Analysis Pipeline
% This reproduces the three-layer pipeline diagram with the fairness-aware contrastive pretraining objective

\begin{figure}[t]
\centering
\caption{Overall Fairness-Aware Histopathology Analysis Pipeline. \textbf{(A)} Three-layer methodology spanning data, training, and evaluation. \textbf{(B)} Fairness-aware contrastive pretraining objective that decouples biological signal from technical and demographic attributes.}
\label{fig:overall-pipeline}
\begin{tikzpicture}[
    node distance=2mm,
    >=latex,
    box/.style={
        rectangle,
        rounded corners,
        draw=black,
        thick,
        align=center,
        font=\small
    },
    layerbox/.style={
        rectangle,
        rounded corners,
        thick,
        align=center,
        font=\small
    },
    titlebox/.style={
        rectangle,
        draw=black,
        thick,
        fill=white,
        align=center,
        font=\bfseries\small
    },
    arrow/.style={
        ->,
        thick,
        >=stealth
    }
]

% Define colors from the original figure
\definecolor{datalayercolor}{RGB}{197,199,237}
\definecolor{traininglayercolor}{RGB}{213,232,212}
\definecolor{evaluationlayercolor}{RGB}{200,206,217}
\definecolor{pretraincolor}{RGB}{255,217,195}
\definecolor{harmonisationcolor}{RGB}{197,216,237}
\definecolor{finetuningcolor}{RGB}{229,241,200}
\definecolor{calibrationcolor}{RGB}{200,220,207}
\definecolor{evaluationboxcolor}{RGB}{195,205,237}

% ==================== MAIN TITLE ====================
\node[font=\bfseries\Large, anchor=south] (main-title) at (0,0) {Overall Fairness-Aware Histopathology Analysis Pipeline};

% ==================== LAYER TITLES ====================
\node[below=3mm of main-title.south west, anchor=north west, font=\bfseries\large] (data-title) {Data Layer};
\node[below=3mm of main-title.south, anchor=north, font=\bfseries\large] (training-title) {Training Layer};
\node[below=3mm of main-title.south east, anchor=north east, font=\bfseries\large] (eval-title) {Evaluation Layer};

% ==================== DATA LAYER ====================
\node[below=1mm of data-title.south, anchor=north, layerbox, fill=datalayercolor, minimum width=3.2cm, minimum height=3.8cm] (data-layer) {};

\node[below=2mm of data-title.south, anchor=north, align=center, font=\small] (data-content) {
    \textbf{Combined Data} \\
    \textbf{Input Layer} \\[2mm]
    \begin{tabular}{@{}l@{}}
    \textbullet~Multi-site WSIs \\
    (20+ sites, multiple \\
    scanners) \\
    \textbullet~Paired clinical slides + \\
    Demographic metadata \\
    (Block C)
    \end{tabular}
};

% ==================== TRAINING LAYER ====================
\node[below=1mm of training-title.south, anchor=north, layerbox, fill=traininglayercolor, minimum width=9.5cm, minimum height=3.8cm] (training-layer) {};

% Training sub-boxes
\node[anchor=north west, layerbox, fill=pretraincolor, minimum width=2.2cm, minimum height=3.6cm] (pretrain) at ([xshift=2mm,yshift=-2mm] training-layer.north west) {
    \textbf{1) Multi-task} \\
    \textbf{adversarial} \\
    \textbf{pretraining} (see B) \\[1mm]
    \begin{tabular}{@{}l@{}}
    \textbullet~Contrastive \\
    learning with multi- \\
    task adversarial \\
    constraints \\
    \textbullet~Forcing technical \\
    site/stain (TSS) and \\
    attribute invariance
    \end{tabular}
};

\node[right=1mm of pretrain.east, anchor=west, layerbox, fill=harmonisationcolor, minimum width=2.2cm, minimum height=3.6cm] (harmonisation) {
    \textbf{2) Multi-attribute} \\
    \textbf{feature} \\
    \textbf{harmonisation} \\[1mm]
    \begin{tabular}{@{}l@{}}
    \textbullet~Harmonising \\
    across sites and \\
    demographic \\
    groups (e.g., age, \\
    race, sex) \\
    \textbullet~Domain-adversarial \\
    neural \\
    harmonisation
    \end{tabular}
};

\node[right=1mm of harmonisation.east, anchor=west, layerbox, fill=finetuningcolor, minimum width=2.2cm, minimum height=3.6cm] (finetuning) {
    \textbf{3) Group-specific} \\
    \textbf{fine-tuning} \\
    \textbf{(FairMIL)} \\[1mm]
    \begin{tabular}{@{}l@{}}
    \textbullet~Multiple Instance \\
    Learning (MIL) for \\
    patient-level \\
    survival/risk \\
    prediction \\
    \textbullet~Group-sensitive \\
    objective function
    \end{tabular}
};

\node[right=1mm of finetuning.east, anchor=west, layerbox, fill=calibrationcolor, minimum width=2.2cm, minimum height=3.6cm] (calibration) {
    \textbf{4) Post-} \\
    \textbf{processing} \\
    \textbf{model calibration} \\[1mm]
    \begin{tabular}{@{}l@{}}
    \textbullet~Post-hoc \\
    subgroup score \\
    calibration \\
    \textbullet~Equalizing \\
    performance \\
    metrics
    \end{tabular}
};

% Vertical label for frozen features
\node[right=1mm of pretrain.east, anchor=south, rotate=90, font=\small] at ([yshift=-1.8cm]pretrain.east) {frozen feature extractor};

% ==================== EVALUATION LAYER ====================
\node[below=1mm of eval-title.south, anchor=north, layerbox, fill=evaluationlayercolor, minimum width=3.5cm, minimum height=3.8cm] (eval-layer) {};

\node[below=2mm of eval-title.south, anchor=north, align=center, font=\small] (eval-content) {
    \textbf{Comprehensive} \\
    \textbf{Evaluation Checklist} \\[2mm]
    \begin{tabular}{@{}l@{}}
    \textbullet~Preserved-site CV -- \\
    internal validation \\
    \textbullet~Demographic subgroup \\
    performance \\
    (Age, Race, Sex) \\
    \textbullet~TSS leakage probe -- \\
    site/scanner \\
    classification
    \end{tabular}
};

% ==================== ARROWS BETWEEN LAYERS ====================
\draw[arrow] (data-layer.east) -- ([xshift=-2mm]pretrain.west);
\draw[arrow] (calibration.east) -- (eval-layer.west);

% ==================== EQUATION BOX ====================
\node[below=5mm of training-layer.south, anchor=north, box, fill=white, minimum width=10cm] (equation-box) {
    \parbox{\linewidth}{
        \centering
        \textbf{Fairness-aware contrastive pretraining objective} \\[3mm]
        \[
        \mathcal{L}_{\text{fair}} = -\sum_{i \in B} \log \left[
        \frac{\exp(\text{sim}(\mathbf{z}_i, \mathbf{z}_i^{+}) / \tau)}
        {\sum\limits_{j \in B, i \neq j} \exp(\text{sim}(\mathbf{z}_i, \mathbf{z}_j) / \tau) + \sum\limits_{\text{attr}} \lambda_{\text{attr}} \exp(\text{sim}(\mathbf{z}_i, \mathbf{z}_j^{+}) / \tau)}
        \right]
        \] \\[2mm]
        \begin{tabular}{ll}
        \hline
        \textbf{Term} & \textbf{Definition} \\
        \hline
        $\mathbf{z}_i$ & Feature embedding for image $i$ \\
        $\mathbf{z}_i^{+}$ & Positive pair (different view of same patient) \\
        $\lambda_{\text{attr}}$ & Adversarial penalty weights for attributes \\
        & (e.g. TSS, race, etc.) \\
        \hline
        \end{tabular} \\[2mm]
        \parbox{\linewidth}{\raggedright
        \textbf{Effect:} Effectively decouples biological signal from technical \\
        (site/stain) and group-related (age, race, sex) attributes.
        }
    }
};

\end{tikzpicture}
\end{figure}
